\scanpage{99}

\uline{Proposition} Supposons \deleted{$k$ parfait et $G$} \inserted{$G$ \uline{connexe} de t.f., et} réductif. Alors $P$ est un sous-groupe \uline{parabolique}, \deleted{\unsure{et}} $U$ son radical unipotent, \inserted{et leur formation commute à \unsure{tte ext. du corps|tt chgt} de base.} Dans tous les cas, $G$ \deleted{\unsure{connexe}} de t.f. $\Longrightarrow U$ unipotent$\ldots$

% slanted
\margin{\unsure{Inutile|Peut-être|Vérifier} que \uline{$G$ soit} \uline{connexe} \deleted{\unsure{pour}} \deleted{\unsure{que}}}

Supposons d'abord $G$ réductif. En car. $0$, l'assertion résulte aisément de la théorie infinitésimale. En car. $p>0$, \deleted{on \unsure{note} que l'on \unsure{est ramené|a des racines|des voisins} \unsure{de noter que} la formation de} il faut sans doute travailler un peu$\ldots$

\uline{Corollaire} Toute $\otimes$-filtration d'un foncteur fibre sur un corps parfait [ou sur un anneau quelc.] splitte $\ldots$

\deleted{\uline{\unsure{Proposition|Démonstration}}}

En effet, $H^1(k,U)=0$ par Rosenlicht$\ldots$

\inserted{\uline{Remarque}} Soit $T=\operatorname{Im} j$, et soit $U'$ un ouvert dense dans $U$. Alors (si $k$ \deleted{\unsure{parfait alg clos}} infini) l'ensemble des $uTu^{-1}$ ($u\in U'$) engendre \deleted{\unsure{le groupe} \illegible{2}} le groupe \deleted{$\unsure{P}$} $T.U$ (produit $1/2$ direct)

$\{$En effet, soit $Q$ le gpe engendré. Alors son algèbre de Lie $\underline Q$ contient les $\operatorname{ad}(u)X$ ($u\in U$), donc les \deleted{\unsure{les}} $[u,x]$, avec $u\in\mathfrak u$, donc [en car. $0$] les $\mathfrak g^\alpha,\alpha>0$ [car $[\mathfrak g^\alpha,x]=\alpha\,\mathfrak g^\alpha=\mathfrak g^\alpha$], donc $\mathfrak u\}$.

% slanted
\margin{Il est facile de donner une \unsure{dém} pas \deleted{\illegible{1}} inf. qui marche en \deleted{toute} car.}

D'autre part, pour les représentations $V$ de $G$, \deleted{\unsure{les}} on a

\[
\deleted{\unsure{V^\alpha}}\ \deleted{\unsure{V^G=}} V^{T.U}\left[=\bigcap_{u\in U'} V^{uTu^{-1}}\right]
\]

\deleted{En effet, on sait que $V^G=V^P$, $P$ étant parabolique, or $V^P=V^{T.U}$, \unsure{comme on voit en} \unsure{se}}

\deleted{en remplaçant $V$ par $V^U$ \unsure{(qui est un|(qui est une} on est ramené à voir que pour tt module $V$ sur $P/U$, on a $V^{P/U}=V^T$. \unsure{Or ceci} Or}

\deleted{si $G$ est de rang $1$, on a $\unsure{L^0}$ semi-simple de rang $1$, on a d'ailleurs}
% archive sheet number: 98
